\subsection{Objetivo 1 --- o Oracle como limite superior em DCS/DES}

\textit{Estudar o conceito de Oracle em MCS, com ênfase em sua utilização como
medida de limite superior em técnicas de DCS e DES.}

O Oracle é o seletor ideal: dada uma amostra, ele escolhe um classificador correto
se existir algum no pool. Por isso é usado como cota superior de qualquer método de
seleção --- nenhum seletor real pode acertar uma amostra que \emph{nenhum} membro do
pool acerta. Souza e Cavalcanti (2017) documentam que essa cota é frouxa: técnicas
DCS selecionam classificadores competentes para no máximo cerca de $85\%$ das
instâncias, mesmo com Oracle de $100\%$ no treino.

O experimento reproduz e quantifica essa frouxidão nas $29$ bases. O \On{1} médio
fica entre $0.9951$ e $0.9976$ nos três modos --- ou seja, \emph{praticamente não
discrimina entre pools}. A informação está na distância até os métodos reais, não no
valor do Oracle. Esta é a motivação empírica direta da generalização proposta pelo
subprojeto e o fio condutor das seções seguintes.
